%%%%%%%%%%%%%%%
%
% $Autor: Wings $
% $Datum: 2020-02-24 14:30:26Z $
% $Pfad: PythonPackages/Contents/General/CloudNativeSQL.tex $
% $Version: 1792 $
%
% !TeX encoding = utf8
% !TeX root = PythonPackages
% !TeX TXS-program:bibliography = txs:///bibtex
%
%
%%%%%%%%%%%%%%%


% iX 1/2023 S. 42
% source: https://www.heise.de/select/ix/2023/1/2208108261315352338


\chapter{Fünf Cloud-native SQL-Datenbanken}

\section{Einführung}


Auch in modernen Cloud-ready-Anwendungen geht ohne Datenbanken fast nichts. Hier mischen Cloud-native Datenbanken gerade den Markt auf. iX stellt fünf Vertreter mit SQL-Schnittstelle vor: YugabyteDB, Citus, TiDB, Vitess und CockroachDB.

\begin{itemize}
  \item Datenbanken spielen als zentraler persistenter Speicher für das Gros der Nutzdaten auch in modernen Cloud-Set-ups eine wichtige Rolle.
  \item In der Cloud sind Datenbanken jedoch mit anderen Anforderungen konfrontiert und müssen mehr Funktionen bieten: Skalierbarkeit etwa ist ebenso wichtig wie die Fähigkeit, auf viele Knoten einer Datenbank parallel schreibend zuzugreifen.
  \item Am Markt etabliert sich derzeit mit den Cloud-nativen Datenbanken ein neuer Datenbanktyp: Die behaupten, bereits ab Werk an die Voraussetzungen der Cloud angepasst zu sein und daher in der Cloud besonders gut zu funktionieren.
\end{itemize}

Datenbanken sind ein fester Bestandteil moderner Cloud-Anwendungen. Derzeit schicken sich die Cloud-nativen Datenbanken an, den Platzhirschen MySQL und Co. Konkurrenz zu machen. Doch was ist eine solche Cloud-native Datenbank überhaupt, und welche Werkzeuge stehen zur Verfügung?

State ist ein Zauberwort, mit dem sich Sysadmins in eine lebhafte Diskussion verwickeln lassen. Stateful und Stateless sind dabei zwei Begriffe, die sich früher üblicherweise auf die Art und Weise bezogen, wie Anwendungen Netzwerkverbindungen handhaben. Heute kommt der Begriff State aber regelmäßig auch im Kontext von Daten vor, mit denen Anwendungen hantieren. Und wie üblich sind sich Administratoren und Entwickler keinesfalls einig, ob State bei Anwendungen grundsätzlich eine gute Sache oder eine Ausgeburt der Hölle ist.

Zumindest im Kontext von Clouds scheint das Thema aber durch: Hier gilt heute das Mantra, dass Anwendungen so wenig Daten wie möglich selbst halten sollen, weil jede Instanz jedes Dienstes grundsätzlich als Wegwerfartikel gilt, der sich jederzeit woanders wieder aus dem Boden stampfen lassen muss. In Clouds und Cloud-nativen Set-ups müssen die Daten daher an zentraler Stelle gelagert sein, sodass alle berechtigten Dienste und Komponenten darauf zugreifen können. Üblicherweise realisieren das die meisten Installationen über Datenbanken. Was vielen Entwicklern und Administratoren allerdings im Verlauf der vergangenen Jahre erst langsam klar geworden ist: Das neue Einsatzszenario in der Cloud verändert natürlich auch die Voraussetzungen, die Datenbanken erfüllen müssen.

Früher war es Usus, eine zentrale Datenbank in einem Set-up zu haben, die mit den Mitteln klassischer Hochverfügbarkeit auf Redundanz getrimmt wurde. Wer gerade zitternd an Pacemaker und Co. denkt, weiß zumindest für Linux, wovon die Rede ist. Und der nicht besonders beliebte Cluster-Manager für Linux kam nicht alleine daher, sondern hatte üblicherweise noch Werkzeuge wie DRBD dabei, die die Replikation von Daten im Hintergrund sicherstellen. Skalierbarkeit war nur in die Höhe ein Thema: Weil sich viele Unternehmen beim Design ihrer Datenbanken so wenig Mühe gaben, dass diese absurde Rechenkapazitäten verschlangen, wurden oft Rufe nach neuer, potenterer Hardware laut. Die Last stattdessen mit mehreren Datenbank-Backends abzufedern, hat sich als Prinzip in konventionellen Umgebungen nie etabliert.

In Clouds ist das naturgemäß anders: Hier herrscht die Economy of Scale, ein wirtschaftlicher Betrieb erfordert, dass zu jedem Zeitpunkt jede Komponente eines Set-ups in die Breite erweiterbar ist. Das schafft mehr Bandbreite und nutzt der Performance. Es bewirkt aber auch, dass sich Set-ups sinnvoll skalieren lassen, etwa wenn später eine größere Datenbank als ursprünglich geplant benötigt wird.

\section{Cloud-native Datenbanken existieren}

Am Markt der Datenbanken stiftet das einige Unruhe. PostgreSQL, MySQL, das mit MySQL noch immer weitgehend kompatible MariaDB und andere etablierte Datenbanken würden gern besser in die Cloud passen, scheitern meist aber schon an den technischen Grundvoraussetzungen. Gleichzeitig zieht am Horizont so etwas wie eine neue Generation von Datenbanken auf, die zwar von außen (oft) wie MySQL oder PostgreSQL aussehen, hinter den Kulissen aber völlig anders arbeiten. Hier sind die Voraussetzungen der Cloud in der Architektur der Datenbanken berücksichtigt, sodass diese sich perfekt beispielsweise in Mikroarchitekturanwendungen in Containerumgebungen einfügen. Ihnen allen ist gemein, dass sie nahtlos in die Breite skalieren können, ab Werk keinen Single Point of Failure mehr haben und dass sie die etablierten Datenbanken in Sachen Performance regelmäßig alt aussehen lassen. Grund genug, sich ein paar Vertreter dieser Gattung näher anzusehen: iX stellt hier die fünf Cloud-nativen Datenbanken YugabyteDB, Citus, TiDB, Vitess sowie CockroachDB vor und vergleicht sie miteinander.

Sämtliche Probanden haben oder emulieren SQL-Schnittstellen. Sie folgen dem klassischen Schema der relationalen Datenbanksysteme. Datenbanken mit diesem Protokoll sind nicht die einzigen Cloud-nativen Datenbanken. Parallel dazu entwickelt sich ein ganzer Zweig von Datenbanken, die andere Formate und Protokolle nutzen, etwa Zeitreihendatenbanken oder Key-Value-Speicher. Üblicherweise spricht man zur groben Unterscheidung von SQL- sowie NoSQL-Datenbanken. Dieser Text geht ausschließlich auf Datenbanken ein, die nach außen SQL emulieren. Denn auch wenn sich die Architektur von Datenbanken unter der Haube ändert, sind die Anforderungen der meisten Programme an eine Datenbank heute SQL-Fähigkeit und ein relationales Modell auf Basis von Tabellen, Zeilen und Reihen.

\section{YugabyteDB: Sharding und viel Abstraktion}

YugabyteDB haben wir bereits in einem Artikel in iX 5/2022 vorgestellt [1], trotzdem sollen an dieser Stelle die wichtigsten Details des Designs nicht fehlen. Was bei einem Blick auf sein Architekturdiagramm (Abbildung 1) sofort ins Auge fällt: YugabyteDB exponiert zur Außenwelt hin zwar einen SQL-Dialekt, doch findet sich von einer klassischen relationalen Datenbank unter der Haube keine Spur. Stattdessen setzt das Produkt zum Speichern von Daten auf eine eigene Datenbank namens DocDB. Deren Kern indes basiert auf RocksDB, dem von Facebook konzipierten, extrem schnellen Key-Value-Speicher. Anders formuliert: Auch wenn eine Applikation Daten in Yugabyte über eine SQL-Schnittstelle ablegt, landen diese letztlich als Key-Value-Paar in RocksDB. Dass die Datenbank zur Außenwelt hin nicht nur SQL spricht, sondern auch Schnittstellen für die Protokolle der Key-Value-Speicher Redis und Cassandra bereitstellt, liegt da auf der Hand.


\begin{figure}
	
	\begin{center}
		
		\includegraphics[width=\textwidth]{CloudNativeSQL/CloudNativeSQL01}
		
		\caption[YugabyteDB]{YugabyteDB ist ein Key-Value Store, der mehrere Standardschnittstellen zur Außenwelt exponiert, unter der Haube massiv auf Sharding setzt und somit typisch für die Klasse der skalierbaren Cloud-nativen Datenbanken ist. Quelle: Yugabyte}
	\end{center} 
\end{figure}	



Wie fast alle verteilten Lösungen arbeitet auch YugabyteDB im Hintergrund auf der Ebene der DocDB mit Sharding. Sämtliche Probanden in diesem Test kümmern sich um das Verteilen ihrer Shards komplett autark. Der Administrator greift nur ein, um neuen Speicherplatz zur Verfügung zu stellen, falls der alte ausgeht. Bei YugabyteDB heißen Shards Tablets.

Eingehende Daten verteilt die Software gleichmäßig über alle im Set-up vorhandenen DocDB-Instanzen; dabei ist jede Zeile im klassischen SQL-Denken ein eigener Key-Value-Eintrag in einer DocDB. Für die Verteilung der Shards sind die sogenannten TServer zuständig. Auf jedem Knoten läuft zudem eine Master-Instanz. Sie verwaltet die anfallenden Metadaten und versetzt YugabyteDB in die Lage, aus den Key-Value-Einträgen in den Tablets über sämtliche Datenbanken hinweg wieder zusammenhängende Ergebnisse für die Antwort auf eine SQL-Anfrage zu generieren. Die gesamte Magie bleibt den Clients allerdings verborgen. Die Entwickler geben an, dass ihre Software weitgehend mit dem PostgreSQL-Dialekt kompatibel ist. Die wenigen nicht implementierten Funktionen sind in der Dokumentation aufgelistet. Sie betreffen aber Features, die nur in wenigen Set-ups tatsächlich in Verwendung sein dürften, sodass sich YugabyteDB in den meisten Fällen als Drop-in-Replacement für PostgreSQL nutzen lässt.

\section{Raft für den Konsens}

Wie alle verteilten Systeme steht auch YugabyteDB vor der Herausforderung, einen konsistenten Zustand aller Daten über alle Knoten hinweg sicherzustellen. Darin enthalten sind freilich auch Situationen, in denen ein Cluster wegen Netzwerkproblemen in mehrere Partitionen zerbricht. Die meisten Produkte am Markt setzen hier auf den Konsensalgorithmus Paxos, YugabyteDB nutzt das alternative Raft-Protokoll. Die Datenbank ist implizit redundant, sodass der komplette Ausfall einer Instanz keinen Datenverlust nach sich zieht. Raft repliziert dafür auf der Ebene der TServer seine Shards beziehungsweise Tablets, wobei für jedes Tablet im gesamten Cluster stets ein primärer TServer und beliebig viele sekundäre TServer aktiv sind. Nur die aktive Instanz lässt sich für den Zugriff nutzen.

Weil es in einem Cluster beliebig viele Shards geben kann, bündelt YugabyteDB beim Zugriff auf Daten durch diese Art der Verteilung die Bandbreite aller verfügbaren Backends. In Sachen Performance hat das Produkt daher ausreichend Reserven.

\section{Gute Cloud-Integration}

Technisch präsentiert sich Yugabytes Datenbank auf der Höhe der Zeit, doch stellt sich auch die Frage nach der Integration in Cloud-Umgebungen und verteilte Systeme. Grundsätzlich gilt, dass gerade in eher konventionellen Set-ups nachempfundenen Umgebungen Entwickler mit YugabyteDB alleine kaum glücklich werden. Denn die meisten Clients, die SQL für den Zugriff auf eine Datenbank nutzen, beherrschen schlicht die Funktion nicht, mehrere IP-Adressen für den Server zu verwalten. Ein Loadbalancer ist also nötig, was sich gerade in Cloud-ready-Umgebungen auf Containerbasis etwa mittels Kubernetes aber leicht realisieren lässt. Wer es etwas umfangreicher will, kann auf externe Produkte wie Istio zugreifen, das Clients automatisch auf eines der Datenbank-Frontends leitet.

Da passt es gut ins Konzept, dass YugabyteDB sich gerade in Kubernetes besonders leicht in Betrieb nehmen lässt. Fertige Container mitsamt einer fertigen Integration für Kubernetes existieren, sodass sich YugabyteDB aus Kubernetes heraus wie ein echter First-Class Citizen steuern lässt.

In Summe präsentiert sich YugabyteDB aufgeräumt und komplett: Skalierte und auch skalierbare Datenbanken sind der Kern der Lösung und funktionieren gut und performant. Die Protokolle zur Anbindung an die Außenwelt sind vielfältig und decken SQL- wie Nicht-SQL-Bedürfnisse gut ab. Hinsichtlich der Integration in moderne verteilte Umgebungen und gerade im Hinblick auf Kubernetes gibt es nur wenig zu bemängeln.

\section{Citus: auch PostgreSQL, aber anders}

Der zweite Kandidat im Vergleich ist Citus (Abbildung 2), das sich von YugabyteDB in etlichen Punkten zentral unterscheidet. Um die Motivation hinter und den Werdegang von Citus zu verstehen, ist ein bisschen PostgreSQL-Geschichte unumgänglich. Anders als andere Datenbanksysteme konnte PostgreSQL nämlich über viele Jahre hinweg nicht mit einer integrierten Replikation aufwarten, die den gleichzeitigen Zugriff auf mehrere synchronisierte PostgreSQL-Instanzen ermöglicht hätte. Für MySQL und später MariaDB gibt es eine solche Umsetzung in Form von Galera seit fast 15 Jahren.

Replikationsansätze gibt es für PostgreSQL mehrere, aber keiner davon hat sich als Standard durchgesetzt. Und einige der am Markt verfügbaren Varianten sind Bestandteil gebündelter Produkte oder kostenpflichtig. Das erklärt auch, wieso die Yugabyte-Macher es für einfacher hielten, gleich eine komplett neue Datenbank zu konstruieren und diese „nur“ mit einer PostgreSQL-kompatiblen Schnittstelle zu versehen, statt noch ein weiteres Plug-in für PostgreSQL zu bauen.

\section{Eigentlich angelegt wie ein Plug-in}



\begin{figure}
	
	\begin{center}
		
		\includegraphics[width=\textwidth]{CloudNativeSQL/CloudNativeSQL02}
		
		\caption[Citus]{Citus baut als einziges Produkt im Test keinen Key-Value Store mit Interface für verschiedene Protokolle, sondern ist ein Plug-in für PostgreSQL. Quelle:	Citus}
	\end{center} 
\end{figure}	



Die Citus-Entwickler indes wollten offensichtlich nicht ganz so radikal bei der Wahl ihrer Mittel vorgehen. Die Datenbank ist im Hinblick auf ihre Architektur ein Mittelding zwischen YugabyteDB und klassischen Ansätzen wie Galera für MySQL. Der zentrale Unterschied ist dabei, dass in Citus echtes PostgreSQL zum Einsatz kommt, sodass die Entwickler keine eigene Kompatibilitätsschicht bauen mussten.

Streng genommen handelt es sich bei Citus also gar nicht um eine eigenständige Cloud-native Datenbank. Vielmehr ist Citus eine Erweiterung für PostgreSQL, die Features wie Sharding und die zentrale Zugriffskontrolle für PostgreSQL bereitstellt und daraus eine skalierbare Datenbank konstruiert. Das kann man sich ähnlich wie bei Hadoop vorstellen: Die Worker in Citus speichern eingehende Inhalte ganz regulär in den ihnen zur Seite gestellten lokalen PostgreSQL-Instanzen. Das Sharding erfolgt dabei anhand der in PostgreSQL genutzten Tabellen. Für eine Tabelle beispiel legt Citus dabei einmal fest, welcher Knoten administrativ verantwortlich ist und auf welchen Knoten die Replicas einer Tabelle landen sollen. Denn auch für Citus gilt, dass es die Daten intern repliziert und damit vor Ausfällen schützt.

Das Prinzip funktioniert aber nur, wenn es – und hier ähnelt Citus dem Ansatz von YugabyteDB – eine zentrale Instanz gibt, in deren Metadaten verzeichnet ist, welcher Worker für eine bestimmte Tabelle zuständig ist. Und nicht nur das: PostgreSQL kennt für Clients schlicht keine Funktion, sich Daten von unterschiedlichen Storage-Backends zusammenzusuchen und daraus eine einheitliche Antwort auf eine Datenbankabfrage zu basteln. Hier kommt der Dienst ins Spiel, den Citus Coordinator nennt: Er verwaltet die Metadaten und weiß, welche Daten wo liegen. Ein Client stellt seine Anfrage im SQL-Dialekt von PostgreSQL stets an den Coordinator, der die Daten für den Client transparent zusammensucht und als einheitliche Antwort ausgibt. Gleichzeitig regelt der Coordinator auch, dass von jedem angelegten Shard im Hintergrund Kopien für die Redundanz bereitgestellt werden. Lesezugriffe erfolgen dann parallel von allen verfügbaren Replicas.

\section{Für den Betrieb ist viel Handarbeit nötig}

Citus präsentiert sich nicht als so universell und vollständig wie YugabyteDB. Die Dokumentation beispielsweise verweist etliche Male auf die PostgreSQL-Doku. So zum Beispiel, wenn es um die Redundanz des Coordinator-Knotens geht: Den muss sich ein Citus-Administrator nämlich komplett selbst bauen, idealerweise unter Einsatz der in PostgreSQL vorhandenen Streaming-Replikation. Damit lässt sich eine Hot-Stand-by-Instanz des Coordinators anlegen, die im Falle eines Falles einen ausgefallenen Knoten ersetzen kann. Allerdings – und das ist der Haken – nicht automatisch: Falls das gewünscht ist, müsste der Entwickler doch wieder in Istio mit automatisch wechselnden IP-Adressen und ähnlichem Komplikationen umgehen.

Und das ist nicht das einzige Komfortproblem bei Citus. So gibt es für Citus keine vernünftige Integration in typische skalierte Umgebungen wie Kubernetes. Fertige Container-Images fehlen ebenso wie die Integration in Kubernetes selbst. Das ist für eine Software, die sich explizit an skalierte Umgebungen richtet, eigentlich eine Ungeheuerlichkeit. Der Selbermach-Faktor ist hier also einigermaßen hoch, sodass Citus nur eine bedingte Empfehlung bekommt. Empfehlung vor allem deshalb, weil sich die Software im Test bewiesen hat und gut funktionierte.

\section{Vitess: eigener Unterbau und MySQL zur Abwechslung}

PostgreSQL und MySQL rangieren heute auf Augenhöhe, was Funktionsumfang, Zuverlässigkeit und Bedienbarkeit angeht. Trotzdem entscheidet sich das Gros der Softwareanbieter bei der Wahl der Datenbank für MySQL. Kein Wunder also, dass Datenbankentwickler am Markt auch Potenzial für eine skalierbare Datenbank mit MySQL-Schnittstelle sahen, und Vitess dürfte zweifelsohne die bekannteste Implementierung einer solchen Datenbank sein. Die Software hat sich im produktiven Einsatz hinlänglich bewiesen und trug bei YouTube eine ganze Weile die gesamte Datenbanklast.



\begin{figure}
	
	\begin{center}
		
		\includegraphics[width=\textwidth]{CloudNativeSQL/CloudNativeSQL03}
		
		\caption[Vitess]{Bei Vitess kommen im Backend zwar echte MySQL-Datenbanken zum Einsatz, doch installiert die Software ein Gateway zu den Clients; die echte MySQL-Datenbank sehen Clients also nicht. Quelle: Vitess}
	\end{center} 
\end{figure}	



Architektonisch gleicht Vitess YugabyteDB in vielerlei Hinsicht, wie ein Blick auf die Grafik (Abbildung 3) verdeutlicht. Dreh- und Angelpunkt in Vitess sind die VTTablets. Das sind Dienste, die einem lokal laufenden MySQL zur Seite gestellt werden, dieses aber lediglich als lokalen Datenspeicher nutzen. Anders als bei Citus wird bei Vitess die lokale Datenbank also nicht zum Anwender durchgeschleift. Stattdessen sind die VTTablet-Instanzen dafür zuständig, die einzelnen Shards, die bei Vitess am ehesten mit Tabellen in einer Datenbank vergleichbar sind, über die Grenzen aller Knoten hinweg zu verteilen und zu verwalten. Was die Arbeit etwas verwirrend macht, ist die eigene Nomenklatur, die die Vitess-Entwickler sich ausgedacht haben: Auf oberster Ebene heißt eine Datenbank (im Normalfall) in Vitess Keyspace. Die Benennung soll aufzeigen, dass es bei Vitess so etwas wie logische Datenbanken gibt – also solche, die der Client von außen zwar als normale MySQL-Datenbank sieht, die im Cluster aber so auf keiner der VTTablet-Instanzen liegen. Ähnlich wie Yugabyte arbeitet Vitess mit einem hohen Maß an Abstraktion.

\section{Zentrale Gateway-Instanzen steuern den Zugriff}

Damit das funktioniert, sind weitere Dienste nötig. Von zentraler Bedeutung ist vor allem das Vitess-Gateway VTGate für Anfragen von Clients. Ihm zur Seite steht ein Topology-Dienst, der sämtliche Vitess-eigenen Metadaten verwaltet. Clients reden mit keiner echten MySQL-Datenbank, sondern mit einer Art Nachbau des MySQL-Serverprotokolls, das im VTGate implementiert ist. Dieses nimmt eingehende Anfragen entgegen, sucht die Information aus den Shards heraus, die über die beliebig vielen VTTablet-Instanzen verteilt sind, und liefert dem Client schließlich eine valide, zusammenhängende Antwort im MySQL-Format.

Mittels CLI oder GUI lassen sich dabei ähnlich wie mit dem MySQL-Client Änderungen an den Vitess-Metadaten vornehmen. Wer beispielsweise eine neue Datenbank, also einen neuen Keyspace, anlegen möchte, tut das mit eben diesen Werkzeugen. Die Anzahl der laufenden Instanzen des Gateways sowie auch des Topology-Dienstes ist dabei nicht begrenzt, sie lassen sich clustern, um Hochverfügbarkeit zu erreichen. Apropos HA: Auch Vitess bietet ab Werk eine interne Replikation, die Datenverlust etwa beim Ausfall einzelner Instanzen verhindert.

Positiv fällt auf, dass das Vitess-Gateway Antworten auf Querys in cleverer Art und Weise kompiliert. Parallel eintreffende identische Querys erledigt das Gateway nicht etliche Male, sondern nimmt stattdessen ein Ergebnis und gibt dieses für alle Abfragen heraus. Droht eine große Datenbank-Query den Dienst in die Knie zu zwingen, gibt es in MySQL derzeit keine zuverlässige Möglichkeit, diesen Vorgang während der Ausführung abzubrechen. Vitess erlaubt das jedoch im VTGateway per Kommandozeilenwerkzeug. Und schließlich führt das VTGateway auch ein sinnvolles Connection Pooling ein: Statt die VTTablets im Hintergrund mit Hunderten und Tausenden gleichzeitig offener Verbindungen zu bombardieren, die alle nur sporadischen Traffic erzeugen, führt Vitess Abfragen im Hintergrund durch wenige offene Verbindungen parallel aus. Das spart Speicher auf den Servern mit laufenden Datenbanken und schont die vorhandenen Ressourcen.

Alles in allem macht Vitess im Test einen guten Eindruck. Vorbildlich ist anders als bei Citus und ähnlich wie bei YugabyteDB die gute Integration in skalierbare Umgebungen wie Kubernetes. Zwar lässt Vitess sich in Paketform auch als Add-on zu MySQL auspacken. Die fertigen Container-Images, die der Anbieter per Docker Hub zur Verfügung stellt, machen es aber leicht, einen Cluster schnell und ohne großes Vorwissen in Betrieb zu nehmen. Dabei agiert Vitess in Kubernetes wie Yugabyte als First-Class Citizen und lässt sich vorbildlich aus Kubernetes heraus steuern. So mögen es Admins im Zeitalter von Cloud-ready.

\section{TiDB: fit für OLTP, OLAP, HTAP und ACID}

Wer die Dokumentation der TiDB liest, dem fliegen die Abkürzungen nur so um die Ohren. Was vorrangig etwas damit zu tun hat, dass die Entwickler die Werbetrommel für ihr Produkt heftig rühren. So rühmt man sich, dass TiDB – das übrigens für Titanium DB steht – Transaktionen unter Einhaltung der ACID-Regeln beherrscht und sich mithin besonders für Workloads eignet, die OLTP-Bedingungen erfüllen müssen. OLTP steht dabei für Online Transactional Processing und beschreibt Workloads, die Transaktionen durchführen. OLAP wiederum steht für Online Analytical Processing und meint die Analyse von Big-Data-Mengen, für die TiDB sich ebenfalls bestens eignen soll. Dass Hochverfügbarkeit und implizite Redundanz der Daten auf der Featureliste nicht fehlen, versteht sich da fast schon von allein.




\begin{figure}
	
	\begin{center}
		
		\includegraphics[width=\textwidth]{CloudNativeSQL/CloudNativeSQL04}
		
		\caption[TiDB]{TiDB ist letztlich auch nur ein Key-Value-Speicher, der verschiedene Protokolle exponiert. Bemerkenswert ist, dass es unterschiedliche Storage-Backends gibt, die sich für verschiedene Datentypen besonders gut eignen sollen. Quelle: TiDB}
	\end{center} 
\end{figure}	



Tatsächlich ist TiDB gar nicht so aufregend oder weltbewegend (Abbildung 4), wenn einem die Funktion von Cloud-nativen Datenbanken im Ansatz klar ist. Denn zwischen Vitess, YugabyteDB und TiDB bestehen deutliche Parallelen. So arbeitet auch TiDB mit Sharding, wofür allerdings ein eigener Dienst namens TiKV zum Einsatz kommt, ein Key-Value Store. Wie bei Yugabyte repräsentieren einzelne Querys in TiDB einzelne Nodes im Key-Value Store, und folglich gibt es auch bei TiDB ein Frontend, das die Daten im Hintergrund zusammensucht und sie dann in einem definierten Protokoll ausgibt. Wie die Vitess-Entwickler haben sich die Entwickler von TiDB für MySQL als Backend entschieden. Die Frontend-Komponente heißt TiDB und tritt stets, wie TiKV auch, geclustert auf und ist implizit redundant.

Der Dienst, der in TiDB die Metadaten verwaltet, heißt Placement Driver Server (PD) und fungiert als eine Art Hirn des gesamten Set-ups. Auch die PDs sind zu einem Cluster zusammengeschaltet und erzwingen ein Quorum, sodass eine Netzwerkpartition keine Probleme verursacht. Weil zu TiDB relativ viele Dienste gehören, ist das Set-up allerdings aufwendiger und im Hinblick auf die Wartung etwas komplexer als bei den anderen Kandidaten. Denn je nach Einsatzzweck gesellt sich an die Seite der TiKV-Server ein weiterer Dienst namens TiFlash: Der nutzt intern eine andere Datenstruktur als TiKV und soll besonders dafür geeignet sein, Analysedaten zu speichern und schnell auszugeben.

Die einzelnen TiDB-Komponenten kommunizieren untereinander über eigene APIs. Die TiDB-Instanzen etwa sprechen über ein eigenes Protokoll namens DistSQL miteinander.

Viele Unternehmen migrieren beim Umstieg in die Cloud Daten aus bestehenden lokalen Diensten wie MySQL in eine Cloud-native Datenbank, um die Performance insgesamt zu verbessern. Ein Set-up mit entsprechend vielen Daten stand im Rahmen dieses Tests jedoch nicht zur Verfügung. TiDB veröffentlicht auf der eigenen Website als einziger Anbieter gemessene Zahlen mit nachvollziehbarem Messprotokoll, die einen Eindruck geben von dem, was möglich ist. Angetreten ist dabei eine lokale, einzelne MySQL-Instanz gegen TiDB.

\section{Ordentlicher Bums in der Datenbank}

Die Daten stammen vom asiatischen Logistikkonzern Ninjavan und lesen sich beeindruckend: Bis zu 60-mal schneller wickelt TiDB demnach verschachtelte JOIN-Befehle ab, als es zuvor bei MySQL der Fall war. Im Mittel liegt die Performanceverbesserung „nur“ beim Zehnfachen, aber auch diese Zahl dürfte vielen leidgeplagten Admins lokaler MySQL-Instanzen bereits die Freudentränen in die Augen treiben. TiDB weist hier darauf hin, dass sich eine Migration in die Cloud auch im Hinblick auf die Performance mehr als auszahlen kann.

Darüber hinaus hebt sich TiDB von der Konkurrenz dadurch ab, dass die Firma eigene Werkzeuge zur Verfügung stellt, um bestehende Datenbanken hin zu TiDB zu migrieren. Wie eine solche Migration funktionieren kann, dokumentieren zumindest Vitess und Yugabyte zwar auch, doch greifen beide Firmen dabei auf externe Werkzeuge zurück, mit denen sich möglicherweise nicht jeder Fall abdecken lässt.

Wenn Admins schließlich vor der Wahl stehen, ob sie lieber Vitess oder TiDB nutzen wollen, dürfte ein entscheidender Faktor die Kompatibilität mit MySQL sein. Je exotischer die Kommandos werden, die beim Betrieb der eigenen Datenbank zum Einsatz kommen, desto höher ist die Wahrscheinlichkeit einer Inkompatibilität. Ein Blick in die Kompatibilitätsanmerkungen von Vitess und TiDB offenbart aber schnell, dass bei TiDB die Liste der nicht implementierten Kommandos viel kürzer ist als bei Vitess. Generell hat TiDB also die bessere MySQL-Emulation und verschafft sich gegenüber der Konkurrenz so einen kleinen Vorsprung.

\section{CockroachDB: komplexer Veteran}

Der fünfte und letzte Proband im Test erweist sich als der mit der komplexesten Grundarchitektur: CockroachDB. Die Datenbank schaufelt Daten durch nicht weniger als fünf Layer, um sie letztlich irgendwann auf die Platte zu schreiben. Doch der Funktionsumfang von CockroachDB gleicht dem der anderen Produkte im Test, und von der Komplexität hinter den Kulissen bekommen Administratoren wie Nutzer kaum etwas mit.





\begin{figure}
	
	\begin{center}
		
		\includegraphics[width=\textwidth]{CloudNativeSQL/CloudNativeSQL05}
		
		\caption[CockroachDB]{CockroachDB ist so etwas wie der Veteran unter den verteilten Datenbanken mit Cloud-native-Funktionalität, wie Yugabyte und TiDB im Kern aber ebenfalls ein monolithischer Key-Value-Speicher. Quelle: CockroachDB}
	\end{center} 
\end{figure}	


CockroachDB ist ein Key-Value-Speicher, nutzt also keine bestehende Datenbank. Damit ähnelt es hinsichtlich seiner Struktur YugabyteDB und TiDB. Seine Entwickler haben dafür ein SQL-Frontend entworfen, das das PostgreSQL-Protokoll implementiert – wie üblich mit manchen Einschränkungen, sodass keine volle Kompatibilität zu aktuellen PostgreSQL-Datenbanken besteht. Wie bei den anderen Kandidaten gilt aber, dass die fehlenden Features nur selten zum Einsatz kommen. Obendrein sind die Inkompatibilitäten gut dokumentiert. Administratoren wissen also, worauf sie sich einlassen.

Das Rad erfinden die CockroachDB-Entwickler aber auch nicht neu, was das Thema des Speicherns von Daten auf Datenträgern angeht. Hier kommt stattdessen RocksDB zum Einsatz, das sich zu einer Art Standard im DB-Umfeld entwickelt.

Jede Cockroach-Instanz hat wenigstens einen Store-Prozess, der für das lokale persistente Ablegen von Daten verantwortlich ist. Allerdings treffen die Store-Prozesse keine Entscheidung darüber, welche Daten auf welcher der verfügbaren Store-Instanzen landen. Das tut eine eigene Komponente, der monolithische KV-Speicher, der im Cockroach-Universum ein eigener Dienst ist.

Die Entwickler des Produktes nehmen ihren Mund allerdings ziemlich voll, was Transaktionssicherheit und Konsistenz angeht: CockroachDB soll verteilte Transaktionen ebenso problemlos beherrschen wie die Probleme beseitigen, die der Ansatz der Eventual Consistency mit sich bringt. Dieser führt unter anderem dazu, dass im täglichen Betrieb von Datenbanken Lese- und Schreibblockaden auftreten, wenn auf ein bestimmtes Objekt mehrere Instanzen eines Dienstes gleichzeitig zugreifen wollen. Die CockroachDB-Entwickler begegnen dem Problem mit dem bereits genannten Layer-Kuchen.

\section{Im Layer-Fahrstuhl nach unten}


Den darf man sich wie eine Art Etagenmodell vorstellen, in dem Daten von oben nach unten durchgereicht werden (Abbildung 5). Anfragen und Eingaben machen Clients per SQL-Protokoll. Cockroach übersetzt den SQL-Befehl dann in eine logische Abfrage für den eigenen Key-Value Store, der die Daten sammelt und eine SQL-konforme Antwort an den Client ausgibt. Dabei beachtet er bereits laufende Querys und sich gerade zutragende Veränderungen.

Viele Begrifflichkeiten kommen einem bekannt vor, hat man sich zuvor mit Vitess beschäftigt: Auch hier ist die Rede von Keyspaces, die in etwa das logische Äquivalent zu einer Datenbank in einer klassischen nicht skalierten SQL-Datenbank sind. Zusätzlich nutzt auch Cockroach Sharding, wickelt dieses aber völlig transparent ab und kümmert sich dabei auch gleich um die Replikation der Daten. Dabei ist es möglich, Failure Domains einzurichten, also etwa verschiedene Racks als Ziele zu definieren. Ähnlich wie bei Yugabyte kommt bei CockroachDB zudem der Raft-Algorithmus anstelle des sonst üblichen Paxos zum Einsatz.

\section{Reif für den Cloud-Einsatz}

Grundsätzlich präsentiert sich CockroachDB als durchaus fit für die Cloud, gerade im Kubernetes-Kontext. Hierfür stellt der Anbieter Nutzern einen eigenen K8s-Operator zur Verfügung, der alle zum Betrieb benötigten Komponenten aufs System holt. Ausführliche Anleitungen zum Deployment von CockroachDB etwa in AWS oder GCE finden sich in der Dokumentation des Anbieters, die gut gelungen ist und zu den besseren im Test gehört.

Cockroach weist darauf hin, dass der Operator aktuell keine CockroachDB-Set-ups erstellen kann, die mehrere geografische Regionen überspannen. Das sei zwar technisch machbar, erfordere in K8s jedoch das händische Ausführen der entsprechenden Schritte. Das schmälert die Freude am Produkt aber nicht. Mit CockroachDB kommt man flott zu einer verteilten Datenbank in der Cloud, die hohen Ansprüchen genügt und gut zu administrieren ist.

\section{Fazit: heiter bis wolkig}


Der Vergleich der fünf Probanden zeigt: Wer skalierbare verteilte Datenbanken mit SQL-Schnittstelle in der Cloud betreiben möchte, hat reichlich Optionen. Drei der vorgestellten Produkte emulieren eine PostgreSQL-Schnittstelle. Yugabyte und CockroachDB nutzen dabei ähnliche Architekturen, Citus fällt etwas aus der Reihe. Zugleich schneidet Citus aus mehreren Gründen im Vergleich am schwächsten ab, was nicht zuletzt auf die fast vollständig fehlende Integration in quasi alle modernen Cloud-Technologien zurückzuführen ist. Ein Produkt, das sich in Kubernetes nicht in wenigen Sekunden in Betrieb nehmen lässt, verdient kaum den Stempel Cloud-ready, auch wenn die Software per se die nötigen Funktionen bietet.

Vitess und TiDB exponieren MySQL und ähneln im Hinblick auf ihre Architektur nicht nur einander, sondern auch YugabyteDB und CockroachDB. Allen vieren ist gemein, dass sie für den Einsatz in der Cloud hervorragend vorbereitet sind und Admins mit allen vier Produkten schnell zu brauchbaren Ergebnissen kommen.

Wie üblich beanspruchen die getesteten Produkte ganz unterschiedliche Eigenschaften für sich: TiDB etwa bewirbt das eigene Produkt besonders mit der hohen Performance, die man im realen Set-up und nicht nur in der Teststellung im Labor erreichen soll.

Signifikante Unterschiede bei der Performance waren im Test aber nicht zu erkennen – es sei jedoch auf die beschränkte Aussagefähigkeit hingewiesen, die von unserem kleinen Test-Set-up ausgeht. Dass TiDB und Vitess vergleichbare Performancewerte erreichen, ist nachvollziehbar: Letztlich lösen sie die Probleme verteilter Systeme auf vergleichbare Art und Weise – was für die PostgreSQL-basierten Kandidaten genauso gilt.

Am Ende dürften also Details wie die implementierten SQL-Befehle das Zünglein an der Waage sein. Wie üblich sollten Admins die für sie infrage kommenden Systeme einem direkten Vergleich unterziehen, um das ideale für den eigenen Use Case zu finden. Zumindest die Aufteilung in die beiden SQL-Varianten PostgreSQL und MySQL hilft dabei, den Kandidatenkreis etwas kleiner zu halten.

Der Bedarf für Cloud-native Datenbanken ist jedenfalls da: Applikationen müssen heute fast selbstverständlich skalierbar und mithin auch implizit redundant sein, und anders als früher können sie sich dabei kaum noch etwa auf die Replizierfähigkeiten zwischengelagerter Speicherschichten verlassen. Das klassische MySQL, das früher auf einem replizierten Volume lag und von Knoten A zu Knoten B wanderte, ist in Clouds so nicht mehr sinnvoll umzusetzen. Gut zu wissen, dass valide Alternativen existieren. 


\begin{table}
  \caption{Cloud-native Datenbanken mit SQL-Schnittstelle}

{\tiny 
  \begin{tabular}{l|l|l|l|l|l}	
                  & Citus                   & CockroachDB                 & TiDB                          & Vitess          & YugabyteDB \\ \hline
    Anbieter      & Citus Data              & Cockroach	                  & PingCAP                       & CNCF            & Yugabyte\\ \hline
    URL	          & \URL{www.citusdata.com} & \URL{www.cockroachlabs.com} & \URL{github.com/pingcap/tidb} & \URL{vitess.io} & \URL{www.yugabyte.com} \\ \hline 
    Lizenz        & AGPL 3.0                & Business Source             & Apache 2.0                    & Apache 2.0      & Apache 2.0, \\
                  &                         & License 1.1                 &                               &                 & POLYFORM\\ \hline
   Storage-Engine & lokales                 & KV-Store, Pebble            & KV-Store, RocksDB             & lokales         & KV-Store, DocsDB \\
                  & PostgreSQL              &                             & (TiKV, TiFlash)               & MySQL           & (RocksDB-Ableger) \\ \hline
    Konsensalgorithmus &  -                 & Raft                        & Raft                          & Paxos-Variante  & Raft\\ 
                  &                         &                             &                               & GCS             & \\ \hline
    SQL-Dialekt   & PostgreSQL              & PostgreSQL                  & MySQL 5.7                     & MySQL           &  PostgreSQL\\ \hline 
    weitere Protokolle & -                  & -                           & -                             & -               & Cassandra, \\ 
                  &                         &                             &                               &                 & Redis \\ \hline 
    als Managed   & ok                      & ok                          & ok                            & -               & ok \\
    Service verfügbar &                     &                             &                               &                 & \\ \hline 
    Kubernetes-   & nur händisch            & ok                          & TiDB Operator                 & Vitess Operator & fertige Container \\
    Integration   &                         &                             &                               & for Kubernetes, & \\
                  &                         &                             &                               & fertige Container & \\
  \end{tabular}
}

\end{table}

\section{Quellen}

\begin{itemize}
  \item Martin G. Loschwitz; Die verteilte Cloud-native-Datenbank YugabyteDB; iX 5/2022, S. 62
  \item Links zu den Produktseiten der Anbieter s. ix.de/zj1q
\end{itemize}

\section{Tasks}


\begin{itemize}
	\item Translation
	\item Improvements
	\item Further readings
	\item Presentation
\end{itemize}

